\section{宏观规律存在性的三种情形：精确、近似与失败}
\label{sec:fault-repair}

本节依次给出三种彼此不同的情形：低复杂度精确商已经存在；每个有限规模都没有非平凡精确任务商、但低复杂度近似商随规模改善；以及任何低误差表示都保持微观尺度。三者分别说明精确理论、近似理论和失败证书承担的作用。

\subsection{正例：置换对称故障--修复系统}

对每个 $N\ge2$，令
\[
B_N=X_N=\{0,1\}^N,
\]
其中 $x_i=1$ 表示第 $i$ 个部件故障，并取恒等微观探针。每一步均匀选择一个部件：若它故障，则以概率 $a\in(0,1]$ 修复；若它正常，则以概率 $b\in[0,1]$ 故障；否则保持不变。记相应核为 $K_{N,a,b}$。任务只判断系统是否完全无故障：
\[
g_N(x)=\one_{\{|x|=0\}},
\qquad |x|:=\sum_{i=1}^N x_i.
\]

\begin{theorem}[故障数商的精确闭合与最小性]
\label{thm:fault-repair}
固定 $0<a_-\le1$，令
\[
\cK_N:=\{K_{N,a,b}:a\in[a_-,1],\ b\in[0,1]\}.
\]
故障数映射 $q_N(x)=|x|$ 是整个核族共同的唯一最粗任务保真精确闭合商。其诱导状态空间为 $Y_N=\{0,1,\ldots,N\}$，且对每个 $(a,b)$ 的宏观核满足
\begin{align}
L_{N,a,b}(k,k-1)&=\frac{ak}{N},\nonumber\\
L_{N,a,b}(k,k+1)&=\frac{b(N-k)}{N},\label{eq:fault-repair-L}\\
L_{N,a,b}(k,k)&=1-\frac{ak+b(N-k)}{N},\nonumber
\end{align}
边界外概率取零。因而
\begin{equation}
c^0_{\cK_N}(g_N)=N+1,
\qquad
\mathcal E_{\cK_N,g_N}(N+1)=0,
\qquad
\Gamma_N\le\frac{N+1}{2^N}\longrightarrow0.
\label{eq:fault-repair-three-levels}
\end{equation}
\end{theorem}

\begin{proof}
从故障数 $k$ 出发，下一步减一、加一或不变的概率只依赖 $k$；直接计数得到式~\eqref{eq:fault-repair-L}，故 $q_N$ 对全部参数同时精确闭合。又有 $g_N(x)=\one_{\{q_N(x)=0\}}$，所以任务保真。

最小性只需考察族中的一个成员 $K_{N,a_-,0}$。设另一个对整个核族精确闭合
的任务保真商合并了故障数 $k<\ell$ 的两个状态。精确闭合使相同宏观初态在
任意未来时刻具有相同任务分布。从 $k$ 个故障出发，前 $k$ 步依次选中各故障
部件并全部修复的概率至少为
\[
\frac{k!}{N^k}a_-^k>0;
\]
从 $\ell>k$ 个故障出发，每步至多修复一个部件，$k$ 步内不可能清零，矛盾。
因此不同故障数必须分开，$q_N$ 就是唯一最粗共同精确商，状态数为 $N+1$。
其分区直径为零，而 $|B_N|=2^N$，其余两式随即成立。
\end{proof}

\begin{corollary}[故障--修复族的完整有限规模前沿]
\label{cor:fault-repair-frontier}
对定理~\ref{thm:fault-repair} 的核族和任务，在
$2\le m\le2^N$ 上有
\begin{equation}
\mathcal D_{\cK_N,g_N}(m)=
\begin{cases}
1/N,&2\le m\le N,\\
0,&N+1\le m\le2^N,
\end{cases}
\qquad
\mathcal E_{\cK_N,g_N}(m)=
\begin{cases}
1/(2N),&2\le m\le N,\\
0,&N+1\le m\le2^N.
\end{cases}
\label{eq:fault-repair-frontiers}
\end{equation}
因此，对每个 $\tau\ge0$，
\begin{equation}
c_{\cK_N}^{\tau}(g_N)=
\begin{cases}
N+1,&0\le\tau<1/(2N),\\
2,&\tau\ge1/(2N),
\end{cases}
\qquad
\Gamma_N=\min\left\{\frac1N,\frac{N+1}{2^N}\right\}.
\label{eq:fault-repair-exact-gamma}
\end{equation}
\end{corollary}

\begin{proof}
先取二块任务分区
$\pi_{\mathrm{task}}=\{\{0^N\},B_N\setminus\{0^N\}\}$。
对固定 $K_{N,a,b}$，非零块中一个故障的状态下一步清零概率为 $a/N$，
两个及以上故障的状态下一步清零概率为零；零故障块是单点。因此
\[
D_{K_{N,a,b}}(\pi_{\mathrm{task}})=\frac aN,
\qquad
\eta_{\pi_{\mathrm{task}}}(K_{N,a,b})=\frac a{2N}.
\]
第二个等式来自区间 $[0,a/N]$ 在 Bernoulli 分布线段上的 Chebyshev 中点。
对参数族取上确界，分别得到 $1/N$ 和 $1/(2N)$ 的可行上界。

再取族内的纯修复核 $K_*:=K_{N,1,0}$。定理~\ref{thm:fault-repair}
证明中的未来任务分布论证同样表明：$K_*$ 的任何任务保真精确闭合分区至少
需要 $N+1$ 块。故任何至多 $N$ 块的任务保真分区 $\pi$ 都满足
$D_{K_*}(\pi)>0$。$K_*$ 从任一状态到每个 $\pi$-块的概率均为 $1/N$
的整数倍，所以两个不同的投影行分布之 TV 距离至少为 $1/N$。于是
\[
D_{\cK_N}(\pi)\ge\frac1N,
\qquad
\eta_\pi(\cK_N)\ge\frac12D_{\cK_N}(\pi)\ge\frac1{2N}.
\]
这与二块任务分区的上界相合，得到式~\eqref{eq:fault-repair-frontiers}；
$N+1$ 块的故障数分区给出其余零值。式~\eqref{eq:fault-repair-exact-gamma}
中的最小状态数由前沿的定义直接得到。最后，预算 $2\le m\le N$ 的
$\Gamma_N$ 目标最小值为 $1/N$，因为 $2/2^N\le1/N$；零直径预算中的
最小目标值为 $(N+1)/2^N$。两者取小即得精确 $\Gamma_N$。
\end{proof}

这个例子中的任务只有“零个故障/至少一个故障”两个输出，而任务未来的精确闭合预测恰需 $N+1$ 个状态。最粗宏观状态记录当前故障数，即保持任务未来所需的最少中间信息。

\subsection{仅近似正例：渐近均匀的异质纯修复系统}

上一例已经具有很小的精确商，因而尚未显示近似闭合相对于精确理论的必要性。现在保留相同的微观空间和任务，却用任意小的逐部件异质性破坏全部非平凡精确任务商。

对 $N\ge2$，仍令 $B_N=\{0,1\}^N$ 并取恒等微观探针。记故障部件集合为
$S(x):=\{i:x_i=1\}$。令
\begin{equation}
 \rho_N:=\frac1N,
 \qquad
 a_{N,i}:=\exp\!\bigl(-\rho_N3^{-i}\bigr),
 \qquad 1\le i\le N.
 \label{eq:heterogeneous-repair-rates}
\end{equation}
每一步均匀选择一个部件；若选中的部件 $i$ 故障，则以概率 $a_{N,i}$ 修复，否则系统保持不变。记所得纯修复核为 $\widetilde K_N$，任务仍为
\[
 g_N(x)=\one_{\{|x|=0\}}.
\]
把它加入定理~\ref{thm:fault-repair}的对称核族，记
\begin{equation}
 \cK_N^\sharp:=\cK_N\cup\{\widetilde K_N\}.
 \label{eq:augmented-repair-family}
\end{equation}
各修复率与 $1$ 的差随 $N$ 消失，所以 $\widetilde K_N$ 是族中纯修复核 $K_{N,1,0}$ 的消失扰动；但对每个固定 $N$，这些微小差异仍编码了故障集合。

\begin{theorem}[只有近似压缩的异质纯修复族]
\label{thm:heterogeneous-repair}
对上述扩充系统族，每个有限规模的最粗精确任务商都是离散分区：
\begin{equation}
 c^0_{\cK_N^\sharp}(g_N)=2^N.
 \label{eq:heterogeneous-exact-complexity}
\end{equation}
而且
\begin{equation}
 d_\infty(\widetilde K_N,K_{N,1,0})\le\frac1{3N}.
 \label{eq:heterogeneous-kernel-perturbation}
\end{equation}
另一方面，故障数表示 $q_N(x)=|x|$ 满足
\begin{equation}
 D_{\cK_N^\sharp}(\ker q_N)
 \le \Delta_N
 \le \frac1{3N},
 \qquad
 \Delta_N:=\max_i a_{N,i}-\min_i a_{N,i},
 \label{eq:heterogeneous-diameter-bound}
\end{equation}
若以 $\Gamma_N^\sharp$ 记扩充核族~\eqref{eq:augmented-repair-family}的联合指标，则
\begin{equation}
 \Gamma_N^\sharp
 \le
 \max\left\{\frac{N+1}{2^N},\frac1{3N}\right\}
 \longrightarrow0.
 \label{eq:heterogeneous-gamma-bound}
\end{equation}
所以该系统族满足 LC，尽管任何有限规模都不存在严格的精确任务压缩。
\end{theorem}

\begin{proof}
先证精确最小性。若一个任务保真精确闭合商合并状态 $x,x'$，则从二者出发的任意未来任务分布都相同。设 $|S(x)|=k<\ell=|S(x')|$。从 $x$ 出发，前 $k$ 步依次选中全部故障部件并成功修复的概率为正；从 $x'$ 出发，每步至多修复一个部件，故 $k$ 步后不可能清零。因此不同故障数的状态不能合并。

再设 $S:=S(x)$ 与 $T:=S(x')$ 是两个不同的 $k$ 元集合。恰在 $k$ 步后清零，必须在这 $k$ 步中各选中 $S$ 内部件一次并全部成功修复，故其概率为
\[
 \frac{k!}{N^k}\prod_{i\in S}a_{N,i}
 =\frac{k!}{N^k}
   \exp\!\left(-\rho_N\sum_{i\in S}3^{-i}\right).
\]
不同有限子集具有不同的三进制和 $\sum_{i\in S}3^{-i}$：取两个子集首次不同的最小指标，其对应项严格大于全部后续项之和。因此 $S\ne T$ 时上述清零概率不同，同一故障数的不同状态也不能合并。所有状态均须分开。任何对 $\cK_N^\sharp$ 共同精确闭合的任务商尤其必须对 $\widetilde K_N$ 闭合，故得到式~\eqref{eq:heterogeneous-exact-complexity}。

再比较 $\widetilde K_N$ 与 $K_{N,1,0}$。从故障集合 $S$ 出发，二者在修复部件 $i\in S$ 的转移概率上相差 $(1-a_{N,i})/N$，相应总质量在 $\widetilde K_N$ 中留于原状态。因此
\[
 \|\widetilde K_N(S,\cdot)-K_{N,1,0}(S,\cdot)\|_{\TV}
 =\frac1N\sum_{i\in S}(1-a_{N,i})
 \le 1-e^{-1/(3N)}
 \le\frac1{3N},
\]
逐行取最大值得式~\eqref{eq:heterogeneous-kernel-perturbation}。

现考察故障数表示。若当前故障集合为 $S$、$|S|=k$，则下一故障数为 $k-1$ 的概率是
\[
 p_S:=\frac1N\sum_{i\in S}a_{N,i},
\]
其余概率留在 $k$。对任意两个 $k$ 元集合 $S,T$，相应两行是同一二点空间上的 Bernoulli 分布，因而
\[
 \|p^{q_N}_{\widetilde K_N,S}-p^{q_N}_{\widetilde K_N,T}\|_{\TV}
 =|p_S-p_T|
 \le\frac{k}{N}\Delta_N
 \le\Delta_N.
\]
由式~\eqref{eq:heterogeneous-repair-rates}与 $1-e^{-x}\le x$，
\[
 \Delta_N
 =e^{-3^{-N}/N}-e^{-1/(3N)}
 \le1-e^{-1/(3N)}
 \le\frac1{3N}.
\]
原对称核族在故障数分区上的直径为零，故扩充族的统一直径由新增核控制。这证明式~\eqref{eq:heterogeneous-diameter-bound}。故障数分区有 $N+1$ 个块，而 $|B_N|=2^N$；把它代入 $\Gamma_N^\sharp$ 的定义即得式~\eqref{eq:heterogeneous-gamma-bound}。
\end{proof}

这个例子把精确与近似两种问题严格分开：一个 $d_\infty$ 意义下趋零的异质扰动把精确任务复杂度从 $N+1$ 推到 $2^N$，但故障数表示的状态比例和闭合误差仍同时趋零。

\subsection{失败例：单点标记循环链}

对 $N\ge2$，令 $B_N=\mathbb Z/N\mathbb Z$，动力学为确定性位移
\[
K_N(x,x+1)=1,
\]
任务为 $g_N(x)=\one_{\{x=0\}}$。

\begin{proposition}[循环链的微观尺度下界]
\label{prop:cycle-obstruction}
对任意 $0\le\tau<1/2$，
\[
c_{K_N}^{\tau}(g_N)=N,
\qquad
\Gamma_N=1.
\]
\end{proposition}

\begin{proof}
任意非离散的任务保真分区都合并某两个状态。若每一对被合并状态的后继仍在同一块中，该等价关系沿循环位移传播，最终会把 $0$ 与某个非零状态合并，违反任务保真。因此必有同块状态的后继落入两个不同块；它们的下一块分布是不同的 Dirac 分布，故纤维预测直径为 $1$，闭合缺陷至少为 $1/2$。于是容限小于 $1/2$ 时只能采用有 $N$ 块的离散分区。对 $\Gamma_N$，离散分区的相对状态数为 $1$，而每个非离散候选的直径为 $1$，所以最小值恰为 $1$。详细的分区论证见附录~\ref{app:examples-counterexamples}。
\end{proof}

置换对称动力学使故障数纤维精确一致；渐近均匀的异质动力学破坏每个有限规模的精确商，却只留下趋零的纤维预测分歧；循环位移则持续传递任务标记的位置差异，使每个非离散任务保真分区的预测直径等于 $1$。三者分别实现精确压缩、仅近似压缩和不可渐近压缩。
